\chapter{Arquitectura multiagente propuesta}

\section{Definición del problema}

El problema se formula como una tarea de decisión secuencial sobre una cartera simulada de activos digitales. En cada instante, el sistema recibe observaciones de mercado, estado de cartera y señales derivadas. A partir de esa información debe decidir si conviene comprar, vender, mantener en observación o rechazar una oportunidad.

La dificultad no está únicamente en predecir si un precio subirá. Una decisión operativa debe considerar el beneficio neto después de comisiones, el cambio de moneda, la liquidez disponible, el capital ya bloqueado, el tiempo de desbloqueo por \textit{trade hold} y la exposición acumulada a un mismo activo o plataforma. Por ello, el objetivo del sistema es maximizar el valor total de la cartera simulada bajo restricciones de riesgo y fricción de mercado.

\section{Visión general del sistema}

La arquitectura separa cinco capas principales: adquisición, persistencia, predicción, decisión multiagente y simulación. Esta separación evita que los agentes de aprendizaje mezclen responsabilidades que no les corresponden. Los extractores observan fuentes externas; el modelo supervisado calcula señales probabilísticas; el entorno MARL convierte esas señales en observaciones; y el simulador aplica las consecuencias económicas de cada acción.

\begin{figure}[h]
\centering
\fbox{
\begin{minipage}{0.9\textwidth}
\centering
Fuentes externas $\rightarrow$ extractores Steam/BUFF/OCR/SteamDT $\rightarrow$ histórico de mercado $\rightarrow$ modelo supervisado calibrado $\rightarrow$ entorno PettingZoo $\rightarrow$ Scout, Trader y Portfolio $\rightarrow$ decisión simulada $\rightarrow$ simulador de cartera
\end{minipage}
}
\caption{Flujo general de la arquitectura propuesta.}
\end{figure}

La decisión final siempre queda registrada como simulada. Esto permite evaluar el sistema con métricas de cartera sin exponer capital real durante la investigación.

\begin{algorithm}
\caption{Ciclo general de toma de decisiones}
\begin{algorithmic}[1]
\State Recoger observaciones recientes del mercado.
\State Actualizar el estado de cartera y las variables de riesgo.
\State Generar una observación para cada agente.
\State Obtener acciones de Scout, Trader y Portfolio.
\State Agregar acciones y aplicar restricciones de riesgo.
\State Simular la decisión final sobre la cartera.
\State Calcular recompensa, métricas y trazabilidad del episodio.
\end{algorithmic}
\end{algorithm}

\section{Definición del entorno}

El entorno representa una cartera simulada que evoluciona a partir de datos históricos. Cada episodio contiene una secuencia de observaciones de mercado y un estado interno formado por saldo disponible, posiciones abiertas, posiciones bloqueadas, fechas de desbloqueo y métricas de exposición.

\subsection{Estado y observaciones}

Las observaciones combinan variables de mercado y variables de cartera. Entre las variables de mercado se incluyen precios de Steam y BUFF163 normalizados a EUR, precios originales en su moneda nativa, volumen, liquidez, spread bruto, spread neto, buy orders, edad de la variante y señales temporales como retornos, medias móviles, RSI o desviaciones de ventas. Entre las variables de cartera se incluyen efectivo disponible, capital bloqueado, posiciones abiertas, exposición por activo y estado de desbloqueo.

La salida del modelo supervisado no se interpreta como orden de compra. Se incorpora como una característica adicional de observación: una probabilidad calibrada de que una oportunidad mantenga un spread rentable en el horizonte definido. De esta forma, los agentes pueden utilizar la señal supervisada sin quedar reducidos a copiarla.

\subsection{Espacio de acciones}

El espacio de acciones se mantiene deliberadamente simple. Las acciones principales son comprar, vender, mantener y observar. En la parte de compra o venta puede añadirse un tamaño de posición discreto para limitar el número de combinaciones. Esta simplificación permite entrenar y depurar el entorno antes de introducir acciones continuas o reglas de asignación más complejas.

\subsection{Transición del entorno}

Cuando se ejecuta una compra simulada, el saldo disponible disminuye y se abre una posición. Esa posición queda bloqueada durante el periodo de \textit{trade hold}. Mientras está bloqueada no puede venderse, pero sí afecta a la exposición y al capital inmovilizado. Cuando se ejecuta una venta simulada, el simulador calcula el ingreso neto aplicando comisiones, conversión de moneda y reglas de valoración. El estado de cartera se actualiza después de cada paso.

\section{Agentes del sistema}

La arquitectura distingue entre agentes extractores y agentes MARL. Los extractores no se entrenan mediante aprendizaje por refuerzo y no deciden operaciones. Su trabajo es adquirir datos desde Steam, BUFF163, SteamDT, OCR o CSV/API y transformarlos en observaciones persistidas. Los agentes MARL, en cambio, operan dentro del entorno de simulación.

\subsection{Scout}

Scout se encarga de detectar oportunidades. Su observación está orientada a señales de precio, spread, liquidez y probabilidad supervisada. Su papel no es decidir toda la operación, sino marcar qué activos parecen merecer atención dentro del episodio.

\subsection{Trader}

Trader decide la acción operativa principal: comprar, vender, mantener u observar. Para ello considera la señal de Scout, el estado de mercado, el posible beneficio neto y el tamaño de posición permitido. Su responsabilidad está más cerca de la ejecución de una oportunidad concreta.

\subsection{Portfolio}

Portfolio controla la coherencia global de la cartera. Su observación incorpora saldo, capital bloqueado, posiciones abiertas, exposición máxima y restricciones de riesgo. Este agente puede penalizar o bloquear decisiones que parezcan rentables de forma aislada pero que deterioren el perfil global de la cartera.

\section{Función de recompensa}

La recompensa debe reflejar el objetivo real del sistema: mejorar el valor neto de la cartera sin ignorar riesgo ni fricción. Una recompensa basada solo en comprar barato o en detectar subida de precio sería incompleta, porque no tendría en cuenta si la operación puede cerrarse, si el capital queda bloqueado o si el spread desaparece antes de poder vender.

\subsection{Recompensa individual}

Las señales individuales pueden ayudar a entrenar comportamientos especializados. Scout puede recibir recompensa por identificar oportunidades que más adelante resultan útiles; Trader por ejecutar acciones con beneficio neto; y Portfolio por evitar sobreexposición, operaciones ilíquidas o capital bloqueado excesivo. Estas recompensas no sustituyen al objetivo global, pero pueden guiar el aprendizaje de cada rol.

\subsection{Recompensa cooperativa}

La recompensa cooperativa se vincula al valor total de la cartera simulada. Debe considerar beneficio realizado, valoración de posiciones abiertas, capital bloqueado, drawdown, operaciones rechazadas por riesgo y costes de transacción. En la implementación actual, el entorno MARL ya expone métricas como recompensa media, operaciones, efectivo, beneficio realizado, drawdown, capital bloqueado y posiciones abiertas.

\section{Restricciones de riesgo}

Las restricciones de riesgo se introducen antes de escalar el entrenamiento. Las más importantes son: no usar capital inexistente, limitar exposición por activo o plataforma, rechazar operaciones con liquidez insuficiente, respetar el \textit{trade hold}, controlar el capital bloqueado y separar claramente simulación de operación real. Estas reglas son sencillas, pero evitan que los agentes aprendan políticas imposibles de ejecutar en el mercado real.
